\section{Beschreibung des Unternehmens}
\label{sec:unternehmen}

\subsection{Handelnde Personen}

Vorstandsvorsitzender (President and CEO) ist Heikki Malinen, zuvor Vorstandsvorsitzender
von Outokumpu (2020--2024) und Posti (2012--2019)\Q{GBXXV}{61}. Er \"ubernahm im Oktober 2024
und startete kurz darauf eine konzernweite \enquote{full potential analysis}, aus der das
laufende Leistungsprogramm hervorging\Q{GBXXIV}{91}. Finanzvorst\"andin ist Eeva Sipil\"a\Q{GBXXV}{62}; sie trat am 13.~Februar 2025 aus dem Board of
Directors aus, dessen Pr\"ufungsausschuss sie geleitet hatte, weil sie zur Finanzvorst\"andin
berufen wurde\Q{GBXXV}{85}.

Das Board of Directors hat acht Mitglieder\Q{GBXXV}{57}. Vorsitzender ist seit 2025 Pasi
Laine, von 2014 bis 2024 Vorstandsvorsitzender von Valmet; stellvertretender Vorsitzender
ist John Abbott, bis 2019 Downstream Director und Mitglied des Executive Committee von Royal
Dutch Shell\Q{GBXXV}{58}. Weitere Mitglieder sind Nick Elmslie (fr\"uher BP Petrochemicals),
Anna Hyv\"onen, Just Jansz, Essimari Kairisto, Conrad Keijzer (Vorstandsvorsitzender von
Clariant) und Sari Mannonen\Q{GBXXV}{59}.

Der finnische Staat sitzt nicht im Board, aber er f\"uhrt den Nominierungsausschuss der
Aktion\"are: Vorsitzende ist Maija Strandberg, Generaldirektorin der Abteilung f\"ur
Eigent\"umersteuerung im Amt des Ministerpr\"asidenten; die beiden weiteren Mitglieder kommen
von den Pensionsversicherern Varma und Ilmarinen\Q{GBXXV}{55}. \bk{Die Besetzung des Boards
ist damit faktisch eine Entscheidung des Mehrheitsaktion\"ars, getroffen \"uber den
Ausschuss statt \"uber eigene Mandate. Fachlich sitzen mehrere fr\"uhere Raffinerie- und
Chemiemanager im Board (Shell, BP, LyondellBasell, Clariant); das passt zu einem Unternehmen, dessen Problem der letzten zwei
Jahre die Verl\"asslichkeit der Anlagen war.}

\subsection{Eigent\"umerstruktur}

\begin{table}[htbp]\centering\small
\caption{Aktion\"arsstruktur zum 31.12.2025 nach Gruppen.}\label{tab:eigentuemer}
\begin{tabular}{@{}lr@{}}\toprule
Gruppe & Anteil\\\midrule
Finnischer Staat & \Pz{\StaatAnteil}\\
Ausl\"andische Aktion\"are & \Pz{\AuslandAnteil}\\
Finnische Institutionen & \Pz{\InstitutionenAnteil}\\
Private Haushalte & \Pz{\HaushalteAnteil}\\
\bottomrule\end{tabular}\end{table}

Die Anteile der Gruppen stehen im Gesch\"aftsbericht 2025\Q{GBXXV}{239}; Neste z\"ahlte
zum Jahresende \Aktionaere{} Aktion\"are\Q{GBXXV}{238}. Ausstehend waren zum Jahresende \AktienMio{} Mio.\ Aktien\Q{GBXXV}{146}; die Zahl ist seit
2017 nahezu unver\"andert (Tabelle~\ref{tab:ergebnis}), R\"uckk\"aufe
weist die Reihe nicht aus. Die vier Gruppen summieren sich auf den vollen Bestand.

\textbf{Meldungen der F\"uhrungskr\"afte (Art.~19 MAR), 27.09.2025 bis 27.09.2026.}
Gesucht wurde in den B\"orsenmitteilungen \enquote{Managers' transactions} auf neste.com und
deren Wiedergabe bei Inderes\QW{Neste, Managers' transactions (Einzelmeldungen)}{quelle:dd}{27.09.2026}.

\begin{table}[htbp]\centering\footnotesize
\caption{Gesch\"afte von Board und Management in zw\"olf Monaten. Sekund\"arquelle, nicht
seitengepr\"uft. Zuteilungen sind Verg\"utungsbestandteile ohne Kaufpreis.}\label{tab:dd}
\begin{tabular}{@{}lllrrr@{}}\toprule
Datum & Person, Funktion & Art & St\"uck & Kurs (EUR) & Volumen (EUR)\\\midrule
\tabellenkoerper{data/dd}
\bottomrule\end{tabular}\end{table}

In zw\"olf Monaten gab es \DdKaeufe{} K\"aufe, \DdVerkaeufe{} Verkauf und \DdZuteilungen{}
Zuteilungen aus Verg\"utungspl\"anen. Den einzigen Verkauf t\"atigte ein Board-Mitglied an
der Londoner B\"orse; er hat mit \DdElmslieStueck{} St\"uck zu \je{\DdElmsliePreis}{EUR}
Kleinstumfang. \bk{Aus den Meldungen l\"asst sich kein Urteil der Organe \"uber den Kurs
ablesen, weder in die eine noch in die andere Richtung.}

\subsection{Beurteilung durch die Kapitalm\"arkte}

Die Aktie ist an der Nasdaq Helsinki notiert. Zum \KursDatum{} kostete sie
\je{\Kurs}{EUR}\QW{Yahoo Finance, NESTE.HE}{quelle:kurse}{27.09.2026}; der B\"orsenwert
betr\"agt \Mio{\Boersenwert}{EUR}, die Spanne der letzten zw\"olf Monate reicht
von \je{\Tief}{EUR} bis \je{\Hoch}{EUR}. Seit Jahresende 2025 hat die Aktie
\KursSeitJahresende{} gewonnen, vom Zw\"olfmonatstief das \KursTiefZuHeute-fache. Das
Kurs-Buchwert-Verh\"altnis liegt bei \KBV. Moody's bewertet Neste mit A3 (Ausblick
stabil)\QW{Neste, Credit ratings}{quelle:rating}{27.09.2026}.

\subsection{Analystenabdeckung und Konsens}

Das Unternehmen ver\"offentlicht keine Ratings und keine Kursziele. Die folgenden Werte sind
Sekund\"ardaten\QW{stockanalysis.com, Neste Forecast}{quelle:konsens}{27.09.2026}.

\begin{table}[htbp]\centering\small
\caption{Analystenkonsens, Sekund\"arquelle, abgerufen am 27.09.2026.}\label{tab:konsens}
\begin{tabular}{@{}lr@{}}\toprule
Kennzahl & Wert\\\midrule
Erfasste Analysten & \KonsensAnalysten\\
Konsensurteil & Kaufen\\
Mittleres Kursziel (EUR) & \KonsensZiel\\
H\"ochstes / niedrigstes Ziel (EUR) & \KonsensHoch{} / \KonsensTief\\
Referenzkurs der Seite (EUR) & \KonsensKurs\\
Ergebnis je Aktie 2026 / 2027, Konsens (EUR) & \KonsensEpsXXVI{} / \KonsensEpsXXVII\\
\bottomrule\end{tabular}\end{table}

\textbf{Aktualit\"atsprobe.} Die Seite nennt einen Aufschlag von \Pz{\KonsensAufschlag} auf
\je{\KonsensKurs}{EUR}; daraus folgt dasselbe Kursziel, und der Referenzkurs ist der
Schlusskurs vom \KursDatum{}, der auch hier verwendet wird. Der Konsens ist damit aktuell.
MarketScreener und Investing.com zeigen dieselbe Zahl der Analysten und dasselbe mittlere Ziel
auf den Cent: Das ist ein Datenstrom, dreimal gesehen, keine drei unabh\"angigen Belege.

Seit dem Halbjahresbericht haben mehrere H\"auser ihre Ziele angehoben, darunter Goldman
Sachs (Heraufstufung auf Kaufen, Ziel \je{\ZielGoldman}{EUR}), Berenberg
(\je{\ZielBerenberg}{EUR}) und Nordea (\je{\ZielNordea}{EUR}).

\subsection{Botschaften und L\"ucken}

Die Botschaft des Konsenses liegt weniger im Urteil als in der Ergebnisfolge: Der Konsens
erwartet f\"ur 2027 ein Ergebnis je Aktie von \je{\KonsensEpsXXVII}{EUR} nach
\je{\KonsensEpsXXVI}{EUR} f\"ur 2026, also einen R\"uckgang um \EpsRueckgangXXVII. Auch die Analysten lesen 2026
als Ausnahmejahr.

Nicht erh\"altlich waren die Einzelziele von Barclays und Morgan Stanley (nur die Urteile
Kaufen und Halten sind \"offentlich), eine Konsensangabe zum EBITDA sowie eine
Netto-Leerverkaufsposition: Im Register der finnischen Finanzaufsicht war am 27.09.2026 keine
Position \"uber der Meldeschwelle von \Pz{\ShortSchwelle} verzeichnet\QW{shortregister.com, Finnland}{quelle:short}{27.09.2026}.
Eine Liste der gr\"o{\ss}ten Einzelaktion\"are au{\ss}er dem Staat fehlt im Gesch\"aftsbericht.

\subsection{Wettbewerbsposition}

Neste beschreibt sich als weltweit f\"uhrenden Hersteller von erneuerbarem Diesel und
nachhaltigem Flugkraftstoff\Q{GBXXV}{4}. Belegt werden kann davon Folgendes:

\begin{table}[htbp]\centering\footnotesize
\caption{Behauptete Vorteile und was die Zahlen zeigen.}\label{tab:wettbewerb}
\begin{tabularx}{\linewidth}{@{}p{2.6cm}Xl@{}}\toprule
Behaupteter Vorteil & Beleg & Tr\"agt?\\\midrule
Gr\"o{\ss}e & Nennkapazit\"at erneuerbarer Produkte rund \je{\Kapazitaet}{Mio.\,t} auf drei
Kontinenten\QI{GBXXV}{14}, nach Rotterdam \je{\KapazitaetZiel}{Mio.\,t}\QI{GBXXIV}{8}. Ein
Marktanteil gegen\"uber einem definierten Markt ist nicht ver\"offentlicht. & teilweise\\
Preissetzungsmacht & Die eigene vergleichbare Verkaufsmarge lag 2025 mit
\je{\RpMarge}{USD/t} um \je{\MargeAbstandRef}{USD/t} \emph{unter} der vom Unternehmen selbst
genannten Referenz-Bruttomarge f\"ur erneuerbaren Diesel (\je{\RpReferenzmarge}{USD/t}),
2024 um \je{\MargeAbstandRefVj}{USD/t}\QI{GBXXV}{81}. & nein\\
Rohstoffbasis & Abf\"alle und Reststoffe stellen \Pz{\Abfallanteil} des Einsatzes
(Vorjahr \Pz{\AbfallanteilVj})\QI{GBXXV}{81}. \bk{Abfallbasierte Kraftstoffe sind die, auf
die Quoten und Anrechnungsregeln zielen.} & ja\\
Kapitalrendite & Ziel war langfristig ein vergleichbarer ROACE von
\Pz{\RoaceZielAlt}\QI{GBXXII}{149}; erreicht 2017 bis 2023 in jedem Jahr, verfehlt 2024
und 2025 (Tabelle~\ref{tab:ergebnis}), in den letzten zw\"olf Monaten
\Pz{\RoaceLtm}\QI{HJ}{3}. & zyklisch\\
Vergleich mit Wettbewerbern & Das Gemeinschaftsunternehmen Diamond Green Diesel meldet
EBITDA je Gallone (\je{\DgdQzwei}{USD} im zweiten Quartal 2026), Neste eine Verkaufsmarge
je Tonne. Die Gr\"o{\ss}en sind verschieden abgegrenzt und ohne weitere Annahmen nicht
vergleichbar. & nicht belegbar\\
\bottomrule\end{tabularx}\end{table}

\bk{Die Tabelle zeigt ein Gr\"o{\ss}en- und Rohstoffunternehmen, kein
Preissetzungsunternehmen: In den beiden schwachen Jahren lag die eigene Marge unter der
Marktreferenz. Der Vorteil tr\"agt die Menge, nicht den Preis.}
